\section{Datengrundlage und Anwendungsrahmen}

\subsection{Alurohr als kontrollierter Testdatensatz}

Die Projektarbeit verwendet eine etwa 48 Sekunden lange Videosequenz eines
Alurohrs als kontrollierten Labordatensatz. Das Rohr repräsentiert ein dünnes,
lineares Objekt, dessen Bildfläche sich mit Kameraposition, Blickwinkel und
Auflösung deutlich verändert. Der Datensatz dient nicht als maßstäblicher
Ersatz für ein im Baugraben liegendes Kabel. Er ermöglicht jedoch, die
Datenübergaben und Fehlerfälle einer vollständigen Rekonstruktionskette unter
wiederholbaren Bedingungen zu untersuchen.

Der Segmentierungsprompt lautet \texttt{pipe}. Für die Parameterstudie werden
zwei zeitliche Abtastungen und drei räumliche Auflösungen verwendet:

\begin{table}[H]
\centering
\begin{tabular}{lll}
\toprule
Profil & Auflösung & Zeitliche Abtastung\\
\midrule
720p & $1280\times720$ Pixel & 2 und 5 FPS\\
QHD & $960\times540$ Pixel & 2 und 5 FPS\\
Low & $640\times360$ Pixel & 2 und 5 FPS\\
\bottomrule
\end{tabular}
\caption{Untersuchte räumliche und zeitliche Frameprofile.}
\label{tab:frameprofile}
\end{table}

Bei rund 48 Sekunden Aufnahmedauer entstehen bei 5 FPS 240 Frames und bei
2 FPS 96 Frames. Alle Verarbeitungsstufen eines Matrixlaufs verwenden denselben
Frame-Satz. Dadurch bleiben Bildnamen, Masken, Kameras und feste
Evaluationsansichten innerhalb eines Laufs konsistent.

\subsection{Scope des linearen Einobjektfalls}

Die Centerline-Nachverarbeitung ist bewusst auf ein einzelnes lineares Objekt
begrenzt. Verzweigte Leitungsnetze erfordern eine zusätzliche Trennung echter
Abzweigungen von Skelettartefakten. Gerade auf dünnwandigen und verrauschten
Meshes kann die Skelettierung Medialflächen und zahlreiche kurze Stachel-Äste
erzeugen. Der produktive \texttt{single}-Modus reduziert deshalb die größte
zusammenhängende Skelettkomponente auf ihren Durchmesserpfad. Der
Netzwerkmodus bleibt außerhalb des Projektarbeitsumfangs.

Das primäre analytische Produkt der PA ist die lokale Centerline. Masken,
COLMAP-Modell, STS-Gaussians und Mesh sind notwendige und diagnostisch wichtige
Zwischenprodukte. CSV-, B-Spline- und GeoJSON-Dateien werden ebenfalls erzeugt,
aber wegen der noch fehlenden unabhängigen Transformations- und
Vermessungsreferenz nicht als global genaue Endprodukte bewertet.

\subsection{Referenzlage und Aussagegrenzen}

Für den aktuellen Testdatensatz liegt keine unabhängig eingemessene
GNSS-Referenzkurve vor. Vorhandene GCP-Dateien bilden nur einen synthetischen
beziehungsweise lokalen Test der Dateiverträge. Fehlt eine validierte
4x4-Ähnlichkeitstransformation, kann die Pipeline einen explizit bezeichneten
Translation-Fallback ausführen. Dieser prüft den Exportpfad, aber weder Rotation
und Maßstab noch Lage- und Höhengenauigkeit.

Daraus folgt die zentrale Abgrenzung: Die PA bewertet Funktionsfähigkeit,
Robustheit, Reproduzierbarkeit, objektmaskierte Ansichtsqualität und relative
Unterschiede zwischen Meshrouten. Ein Nachweis der später angestrebten
Toleranz von $\pm10\,\mathrm{cm}$ ist erst mit realen GCP-/GNSS-Messungen,
einer validierten Transformation und geometrischen Vergleichsmetriken möglich.
